\section{Related Work}

This section reviews the theoretical and empirical foundations of automated IT support ticket prioritization. After positioning ticket prioritization within the broader context of ITSM, it outlines NLP-based feature representations and relevant machine learning classifiers. Finally, XAI is introduced as the foundation for evaluating explainability.

\subsection{Ticket Prioritization in IT Service Management}

Support ticket systems provide a central interface between users and the help desk in ITSM by recording technical problems and service requests as tickets \cite{alhawari2021helpdesk,fuchs2022support}. Tickets typically contain textual problem descriptions as well as additional information such as category or priority \cite{fuchs2022support,zangari2023ticket}. The ticket management process comprises activities such as logging, classification, prioritization, assignment, subsequent processing, and closure \cite{alhawari2021helpdesk}.

Prioritization represents a distinct decision task within ticket processing. In an ITSM context, priority is determined based on criteria such as urgency and business impact \cite{razma2026classification}. This decision is relevant because limited resources require tickets to be processed in a prioritized order and important tickets to be addressed early \cite{tian2015priority}. Tickets may also be subject to time-bound service requirements. Delays can affect business processes, while violations of Service Level Agreements (SLAs) may have financial consequences \cite{gupta2012scheduling}.

Ticket automation covers different tasks such as classification, routing, expert finding, and ticket resolution \cite{fuchs2022support,zangari2023ticket}. Within this broader scope, ticket prioritization constitutes a specific classification task in which incoming tickets are assigned to a priority class \cite{lyubinets2018labeling,beckers2009multifacet}.

\subsection{NLP-Based Feature Representation for Ticket Data}

For automated classification, unstructured ticket texts must be transformed into numerical feature representations. Ticket texts may contain short descriptions, domain-specific terminology, abbreviations, or technical artifacts. Real-world IT support data may additionally involve challenges such as multilingual input and class imbalance \cite{lyubinets2018labeling,razma2026classification}. Beyond the ticket text itself, features derived from the subject line, text position, punctuation, attachments, and sender information can provide relevant information \cite{beckers2009multifacet}.

An established representation for text classification is Term Frequency--Inverse Document Frequency (TF-IDF), which represents documents as high-dimensional, sparse vectors based on weighted terms \cite{manning2008introduction}. TF-IDF is also combined with classical machine learning methods in ticket classification \cite{revina2020ticket}. One limitation is that relationships between terms are captured only to a limited extent \cite{bonechi2024support}.

Contextual language models such as BERT, by contrast, incorporate both left and right linguistic context and can be fine-tuned for downstream classification tasks \cite{devlin2019bert}. Recent ticket studies compare classical TF-IDF-based approaches with transformer-based models \cite{bonechi2024support,razma2026classification}. \citet{bonechi2024support} show that BERT can achieve higher classification performance, whereas TF-IDF combined with Support Vector Machine (SVM) provides a simpler and more resource-efficient alternative.

\subsection{Machine Learning Classifiers for Ticket Prioritization}

Ticket classification has been investigated using different model families, including probabilistic, linear, instance-based, and tree-based methods as well as ensemble, deep learning, and transformer models \cite{revina2020ticket,ahmed2023incident}. Comparative studies indicate that model performance varies depending on the underlying data and classification task, and no uniformly superior model family can be identified across different application contexts \cite{revina2020ticket,ahmed2023incident}. Their respective strengths and limitations for text-based ticket data are therefore particularly relevant for their methodological positioning.

Naive Bayes methods constitute an established probabilistic family for text classification and can be applied efficiently to high-dimensional text data due to their comparatively low computational complexity. Their modeling, however, relies on strong conditional-independence assumptions between features. Limitations may arise particularly in the presence of imbalanced text classes \cite{kibriya2004multinomial,frank2006naive}. Complement Naive Bayes is also used as a comparison method in ticket classification and therefore extends the model comparison by a probabilistic model family \cite{zicari2022ticket}.

For large and sparse feature spaces such as those produced by TF-IDF, Logistic Regression and linear SVMs are particularly suitable, as they learn linear decision functions  \cite{joachims1998text,fan2008liblinear}. This efficiency comes with a limited ability to directly model nonlinear interactions. Nevertheless, recent help desk studies report competitive SVM performance compared with Random Forest, Long Short-Term Memory (LSTM) models, and BERT, in some cases with lower computational resource requirements \cite{martins2024port,bonechi2024support}.

Instead of learning a global decision function, k-Nearest Neighbors (kNN) serves as an instance-based alternative that classifies new instances based on similar training observations \cite{cover1967nearest}. Its performance depends directly on the selected similarity metric and the structure of the feature space. kNN is also used as a comparison method in ticket classification \cite{revina2020ticket}.

Tree-based ensemble methods can model nonlinear relationships and feature interactions. Random Forest aggregates randomized decision trees and thereby reduces dependence on individual tree structures \cite{breiman2001random}. Boosting methods, by contrast, construct models sequentially. XGBoost combines this principle with, among other techniques, regularization and sparsity-aware processing \cite{chen2016xgboost}, whereas LightGBM is specifically designed for efficient gradient boosting on large and high-dimensional datasets \cite{ke2017lightgbm}. This increased model flexibility is accompanied by greater model and tuning complexity.

Recent empirical findings further illustrate the dependence of model performance on the specific application context. In the industrial IT incident data analyzed by \citet{ahmed2023incident}, XGBoost is among the strongest methods on one dataset, whereas convolutional neural networks or transformer models achieve the best results on other subsets. In contrast, \citet{martins2024port} report the highest F1-scores for SVM on real-world help desk requests and substantially shorter training times for classical methods compared with BERT and LSTM. These differing findings provide the rationale for RQ1 and the controlled comparison of multiple model families under a common data, feature, and evaluation basis.

\subsection{Explainable AI and Explanation Quality}

In addition to classification performance, it is relevant to understand the information on which model predictions are based. XAI encompasses approaches intended to make machine learning predictions and model behavior more understandable \cite{guidotti2018survey,arrieta2020xai, adadi2018peeking}. Explainability has also been investigated in the ticket context. \citet{zicari2022ticket} combine ticket classification with local explanation methods to analyze relevant features and misclassifications, whereas \citet{licina2026helpagent} integrate SHAP-based explanations into a help desk system.

Established local explanation methods include Local Interpretable Model-Agnostic Explanations (LIME), which approximates a model's behavior around an individual prediction using an interpretable local surrogate model \cite{ribeiro2016lime}. SHAP, in turn, provides an approach to feature attribution by assigning individual feature contributions to a specific model prediction \cite{lundberg2017shap}. This enables models from different families to be examined based on their attribution patterns. However, explanation quality is multidimensional. \citet{nauta2023evaluation} distinguish twelve properties of explanation quality and include \emph{Compactness} as a separate property concerned with how extensive an
explanation is.

The existing literature therefore includes both diverse approaches to ticket classification and applications of XAI in the ticket context. Although model comparisons focus primarily on classification performance and explanation-oriented studies examine individual XAI applications, cross-model quantitative comparisons of SHAP-based explanation properties in the context of priority prediction remain limited in the reviewed literature. To address this gap, the present study extends the performance comparison in RQ1 with a quantitative assessment of explanation compactness in RQ2.